\documentclass[a4paper,10pt,titlepage]{article}
\usepackage{graphicx}
\usepackage[spanish]{babel}
\usepackage{amsmath}
\usepackage{amssymb}
\usepackage{listings}
\usepackage{multicol}

\title{Planificador de tareas para un sistema hidrop\'onico automatizado programado en C++}
\author{Mauricio Yamil Tame Soria}
\begin{document}
\maketitle
\tableofcontents
\newpage

\section{Introducci\'on}
Se tiene un sistema hidrop\'onico donde se busca automatizar los procesos que mantengan las condiciones \'optimas para el cultivo. Los procesos consisten en regular, 
mediante actuadores, las 
variables del entorno: temperatura ambiente, temperatura en la soluci\'on de riego, humedad relativa en el ambiente, humedad relativa en la tierra y el ph de la soluci\'on 
de riego. Para mantener controladas las variables, es necesario el monitoreo de las mismas mediante muestreos. Dichas variables son de car\'acter 
cuantitativo, cuando se 
haga referencia a ellas, ser\'an n\'umeros 
reales (enteros con decimal). El monitoreo consiste en hacer muestreos de las variables peri\'odicamente. Los valores 
monitoreados, igual que toda la informaci\'on del sistema, son registrados en una base de datos. En los cultivos hidrop\'onicos, la planta se nutre directamente del 
riego por lo que es importante mantener h\'umeda la ra\'iz; se proponen dos formas de riego para cumplir esta condici\'on. La automatizaci\'on del sistema se lleva acabo 
con el planificador de tareas, 
tambi\'en nos referiremos a \'el como ''scheduller''. El planificador de tareas es un programa que mantiene los vegetales regados y procesa las variables del entorno para 
llevar acabo los procesos que 
las mantengan en el rango ideal. Cada cultivo tiene un rango ideal de ph, temperatura ambiente, humedad ambiente y temperatura en la soluci\'on de riego; todos 
ellos se encuentran definidos 
en la 
base de datos. El sistema hidrop\'onico est\'a montado en un microprocesador ''Pcduino''; los componentes (sensores y actuadores) del sistema est\'an 
conectados a los pines de la tarjeta. El planificador de tareas, al igual que todos los programas, fue programado y compilado en el microprocesador; los c\'odigos se 
encuentran en la carpeta principal del proyecto:

\begin{verbatim}
/home/Hidroponia
\end{verbatim}

A lo largo del 
texto se tratar\'an los siguientes temas: configuraci\'on del microprocesador, m\'etodos de muestreo, procesos para regular las variables del entorno, la base de datos 
del sistema, la estructura y 
funcionamiento del 
''scheduller'', pruebas de hardware, los problemas del multithreading y sus soluciones, administraci\'on de los sistemas, estad\'isticas con las variables de entorno, una 
interfaz de usuario y las \'areas de mejora del proyecto.

\section{Configuraci\'on del Pcduino}
La configuraci\'on del Pcduino consiste en la instalaci\'on de un sistema operativo, en este caso una distribuci\'on de Linux (lubuntu), y un conjunto 
de programas que nos permiten montar el servidor y administrar la base de datos. Para la instalci\'on del sistema operativo lubuntu se deben obtener los archivos del 
kernel y el sistema operativo; se encuentran descargables 
en la 
p\'agina del distribuidor del Pcduino: 
\begin{verbatim}
 http://www.linksprite.com 
\end{verbatim}
Una vez descargados los archivos, se debe grabar el archivo imagen del kernel en una microSD con la ayuda 
del comando 
\verb|dd| que provee Linux y guardar el archivo \verb|ubuntu.iso| en una 
USB. Cuando est\'e la imagen del kernel grabada en la microSD, con el Pcduino apagado, se inserta la miscroSD y se conecta el Pcduino a la corriente. Uno de los leds 
parpadea mientras se realiza la instalaci\'on del kernel; cuando deje de parpadear el led se debe retirar la microSD y presionar el bot\'on ''reset'' del duino. Al 
reiniciar el duino, se observa una pantalla pidiendo insertar una USB con la imagen del sistema; se inserta la USB con el archivo \verb|ubuntu.iso| y la instalaci\'on 
del sistema operativo comenzar\'a a ejecutarse. Cuando termine el proceso se mostrar\'a un mensaje en la pantalla, se retira la USB y se reinicia el duino para 
finalizar. Una vez instalado el sistema operativo, se necesita establecer una 
contrase\~na; para esto se abre la l\'inea de comandos (Terminal) y se teclea:
\begin{verbatim}
sudo passwd
\end{verbatim}
ya establecida la contrase\~na se reinicia el Pcduino. Establecer la contrase\~na es importante pues de lo contrario se presentan problemas al momento de 
instalar algunos paquetes. Para configurar la zona horaria, el tama\~no de la pantalla y otras 
preferencias del sistema se ejecuta el comando:

\begin{verbatim}
 sudo board-config.sh 
\end{verbatim}

\subsection{Paquetes}
Los programas necesarios 
para el funcionamiento del servidor son Apache, Mysql para la base de datos, 
php para la aplicaci\'on web (interfaz de usuario) y phpmyadmin para auxiliarnos con la base de 
datos. 
Todos los paquetes se instalaron desde la terminal usando el comando:
\begin{verbatim}
 sudo apt-get install nombre_del_paquete
\end{verbatim}
los nombres de los paquetes son:
\begin{itemize}
\item apache2
\item php5
\item libapache2-mod-php5
\item mysql-server 
\item mysql-client
\item phpmyadmin 
\item libmysqlcppconn-dev 
\end{itemize}

Antes de instalar los paquetes, se debe ejecutar el comando \verb|sudo apt-get update|. Una vez instalados los paquetes, es necesario reiniciar el duino. Al reiniciar, se 
nota que el servidor no 
funciona; el error se debe a que no se crea la carpeta \verb|/var/log/apache2| al momento 
de iniciar el sistema operativo. Una forma sencilla de reparar el ''bug'', es editar el archivo \verb|/etc/rc.local| 
agregando las siguientes l\'ineas de c\'odigo:

\begin{verbatim}
 mkdir /var/log/apache2
 apachectl restart
\end{verbatim}

La carpeta que contiene el 
phpmyadmin se encuentra en \verb|/usr/share/phpmyadmin|; para acceder a ella del navegador, se debe 
crear un link simb\'olico dentro de la carpeta del servidor. Para ello, se ejecuta la 
siguiente l\'inea 
de c\'odigo:
\begin{verbatim}
 sudo ln -s /usr/share/phpmyadmin /var/www
\end{verbatim}

despu\'es de esto ser\'a posible acceder a la aplicaci\'on \verb|phpmyadmin| desde el navegador con la ruta

\begin{verbatim}
http://localhost/phpmyadmin
\end{verbatim}

Una vez instalados todos los paquetes, es necesario reiniciar el sistema. Tras seguir los pasos anteriores el Pcduino estar\'a configurado para el sistema hidrop\'onico.

\subsection{Configuraci\'on de red}
El servidor necesita una ip fija; para esto se debe abrir el men\'u de inicio, pasar a 
preferencias y luego pasar a las conexiones de red donde se edita la conexi\'on que se utilizar\'a. En este caso, se utiliza la conexi\'on por cable ethernet. 
Para fijar la direcci\'on ip se debe pasar a 
la pesta\~na \verb|ipv4| (internet protocol 
version 4) y registrar los datos correspondientes. Actualmente el Pcdino tiene la siguiente configuraci\'on de red:

\begin{itemize}
	\item ip: 10.10.254.78
	\item Netmask: 255.255.0.0
	\item Gateway: 10.10.255.254
	\item DNS server: 8.8.8.8
	\item Search domain: iies.unam.mx
\end{itemize}

Se debe reiniciar el Pcduino para que tomen efecto las configuraciones.

\section{Drivers}
Los drivers son los programas que permiten la comunicaci\'on del sistema con los sensores y 
actuadores. Los drivers son funciones definidas en archivos cpp y declaradas en archivos hpp. El formato 
que se escogi\'o para declarar
funciones de lectura fue el siguiente: 

\begin{verbatim}
 float driver_sensor(int pin);
\end{verbatim}

donde el par\'ametro pin se refiere al puerto donde est\'a conectado el 
sensor. Para el sensor de temperatura en agua, la funci\'on es diferente; 
el par\'ametro de entrada es el identificador del sensor y todos los 
sensores van conectados en paralelo al puerto GPIO 3. GPIO (General Purpose Input/Output, Entrada/Salida de prop\'osito General) es un pin gen\'erico en un chip, 
cuyo 
comportamiento (incluyendo si es un pin
de entrada o salida) se puede controlar (programar) por el usuario en tiempo de ejecuci\'on. Para los 
actuadores, se utilizaron 3 funciones:

\begin{verbatim}
void actuar(int seg,int pin,string mensaje)
void activar(int pin,int pinData,int pinMode)
void parar(int pinData,int pinMode)
\end{verbatim}

La funci\'on actuar recibe los siguientes par\'ametros: \verb|seg| son los segundos que debe funcionar el actuador, \verb|pin| es el puerto donde est\'a conectado el 
actuador y \verb|mensaje| es un texto que sirve para llevar un registro de lo que se ejecuta. La funci\'on activar sirve 
para encender el actuador por un tiempo indefinido y recibe los siguientes par\'ametros: \verb|pin| el puerto donde est\'a conectado el actuador, 
\verb|pinData| es el identificador del archivo donde se almacena la informaci\'on del puerto y \verb|pinMode| es el identificador del archivo donde se define si el puerto 
es de entrada o salida. Por \'ultimo, la funci\'on \verb|parar| es la contraparte de \verb|activar|, es decir, detiene un actuador que se encuentre activado. Para mayor 
detalle, todos los drivers se encuentran 
en 
la carpeta
\begin{verbatim}
 Hidroponia/Codes/Drivers
\end{verbatim}

excepto el driver de temperatura en l\'iquidos, \'este se encuentra en la 
direcci\'on

\begin{verbatim}
Support/c_environment/sample/temagua.c
\end{verbatim}

dentro de la carpeta del proyecto.

\subsection{Protocolo para a\~nadir drivers}
El protocolo consiste en agregar dos archivos a la carpeta

\begin{verbatim}
Hidroponia/Codes/Drivers
\end{verbatim}

un archivo \verb|driver.hpp| que contenga la declaraci\'on de las funciones con el formato anterior y otro archivo \verb|driver.cpp| donde est\'en definidas las funciones 
del driver. Una vez agregados los archivos, se compila el driver ejecutando el binario \verb|comp| dentro de la carpeta \verb|Hidroponia|. El binario \verb|comp| se genera 
del c\'odigo \verb|compiler.cpp|, ubicado en la carpeta Codes. Para crear el binario, es necesario ejecutar dentro de la carpeta 
\verb|Hidroponia| el siguiente comando:

\begin{verbatim}
make compiler
\end{verbatim}

El binario \verb|comp| compila 
todos los archivos \verb|driver.cpp| que se encuentren en la carpeta \verb|Drivers| y a partir de ellos genera los archivos objeto \verb|driver.o|; estos 
son enviados a la carpeta \verb|Hidroponia/Objects|. Una vez que est\'an hechos los 
objetos, el programa \verb|comp| genera un archivo de librer\'ias 
llamado \verb|libs.hpp| 
donde se escriben todos los nombres de los archivos \verb|driver.hpp| que se encuentren en \verb|Drivers|, es decir, se genera un archivo que contiene los headers de todos 
los drivers. El archivo 
\verb|libs.hpp|, es el 
que se importa desde los otros c\'odigos para poder usar las funciones de los drivers. Cuando termina de ejecutarse \verb|comp|, se tienen los objetos que proveen las 
funciones de los drivers y un 
archivo \verb|libs.hpp| para usarlo desde las otras partes del proyecto. El protocolo se pretende automatizar usando el administrador de sistemas; sin embargo la 
automatizaci\'on continua en 
desarrollo y m\'as adelante se expondr\'a tanto el progreso alcanzado como los pasos para finalizarlo.

\subsection{Sensores de temperatura en l\'iquidos}
Los sensores de temperatura en l\'iquidos son controlados por un driver de 
c\'odigo libre que se encuentra en Github; las instrucciones para su 
instalaci\'on est\'an en la p\'agina del 
distribuidor del duino (linksprite). Hay que descargar una carpeta y 
ejecutar el comando \verb|make| que est\'a definido en el archivo \verb|makefile| dentro de la 
carpeta que se descarg\'o previamente. El nombre de la carpeta es \verb|c_environment| y se coloc\'o en la ruta

\begin{verbatim}
 Hidroponia/Support/c_environment
\end{verbatim}

En la Carpeta Support se colocan las herramientas que est\'en programadas en otro lenguaje; este driver est\'a programado en  
$C$. Dentro de la carpeta 

\begin{verbatim}
c_environment/sample
\end{verbatim}

se encuentran algunos drivers para 
diferentes sensores arduino. Para este driver se utiliza el 
archivo llamado

\begin{verbatim}
DS18x20_Temperature.c
\end{verbatim}

Dicho archivo genera un programa que captura lecturas de m\'ultiples sensores conectados en 
un s\'olo puerto GPIO de manera continua y casi aleatoria. El 
programa regresa varios datos para cada lectura: el identificador del 
sensor, un valor en hexadecimal para verificar que la lectura sea 
correcta y la lectura del sensor en celcius. Este driver no permite ejecutar lecturas de un sensor determinado, por lo que fue 
necesario editar el c\'odigo. Se cre\'o una copia llamada

\begin{verbatim}
temagua.c
\end{verbatim}

para realizar sobre ella la edici\'on del c\'odigo y conservar el original. El c\'odigo se modific\'o para que se tomara s\'olo una lectura 
confiable checando el hexadecimal de verificaci\'on; dicha lectura se guarda en un archivo 
junto con el identificador del sensor que la realiz\'o. Para 
crear los binarios de los drivers encontrados en la carpeta \verb|c_environment/sample|, se debe 
ejecutar dentro de la carpeta \verb|c_environment| el comando

\begin{verbatim}
make
\end{verbatim}

para ejecutar el script en el archivo \verb|Makefile| situado en la carpeta. El script compila las 
librer\'ias \verb|Arduino| necesarias para los drivers. Una vez que termina el 
proceso de compilaci\'on de las librer\'ias, se debe pasar a la carpeta \verb|c_environment/sample| para editar el archivo \verb|Makefile| a\~nadiendo la siguiente l\'inea 
de c\'odigo justo despu\'es de la \'ultima l\'inea an\'aloga:

\begin{verbatim}
OBJS += temagua
\end{verbatim}

Esta l\'inea de c\'odigo representa el comando para que se compile el archivo \verb|temagua.c| y se genere el ejecutable \verb|temagua|. Una vez editado el archivo 
\verb|Makefile|, dentro de \verb|sample| se teclea el comando \verb|make| para ejecutar el script y crear el binario el cual es enviado a la carpeta 
\verb|c_environment/output/test|. Para usar el driver 
en el planificador de tareas, se cre\'o una funci\'on en $C++$ que obtiene la lectura de un determinado sensor. El pseudoc\'odigo de la 
funci\'on es el siguiente:

\begin{verbatim}
Temperatura_en_liquido(sensor.id){
  ejecutar temagua;
  leer archivo;
  while(lectura.id!=sensor.id){
    ejecutar temagua;
    leer archivo;
    }
  return lectura;
  }fin driver
\end{verbatim}

Esta funci\'on representa el driver final, es decir, el que ser\'a utilizado por los dem\'as programas. La funci\'on devuelve la lectura en tipo \verb|float| y el 
par\'ametro 
\verb|sensor.id| 
corresponde al identificador del sensor con el cual se efectuar\'a la 
lectura. La declaraci\'on y definici\'on de la funci\'on se hizo en la carpeta \verb|Drivers| sobre los archivos 
\verb|temagua.hpp| y \verb|temagua.cpp| respectivamente.

\subsection{Sensor ph}
El sensor de ph es un sensor de conversi\'on anal\'ogica-digital (adc en 
ingl\'es). Para la calibraci\'on, se tomaron 3 valores conocidos de ph y se ejecut\'o la lectura del sensor sin el driver para cada uno de ellos, obteniendo as\'i 3 puntos 
experimentales. Cada punto experimental representa un par ordenado de voltaje medido por el sensor y el ph de la muestra. Una vez obtenidos los 3 puntos 
experimentales, 
se realiz\'o un ajuste lineal para obtener una recta que 
representa el ph en funci\'on del voltaje medido por el sensor adc. La funci\'on 
qued\'o de la siguiente manera:

\begin{equation}
	ph(V)=V \frac{14*0.8}{4096} 
	\label{ph}
\end{equation}

donde $V$ representa el voltaje registrado 
en el archivo correspondiente al puerto.

\subsection{Temperatura y Humedad ambiente}
La temperatura y la humedad ambiente las provee un s\'olo sensor, el DHT22 de arduino. Como el sensor es digital, no se utilizan los puertos adc sino los GPIO. Lo que hace 
el driver es fijar el pin en modo output para despu\'es 
capturar los bytes y procesarlos. El procesamiento de bytes devuelve lecturas de temperatura en celcius y humedad relativa en porcentaje (de 100). Aunque es un s\'olo 
driver, se crearon dos funciones: 
una para obtener la temperatura y otra para obtener la humedad. Las funciones se encuentran definidas y declaradas en los archivos \verb|temhum.cpp| y \verb|temhum.hpp|. La 
separaci\'on de las funciones tiene el 
fin de seguir el protocolo de funcionamiento de los drivers.

\subsection{Humedad en tierra}
El sensor de humedad en tierra es parecido al de ph, es del formato adc. El sensor consiste en 2 placas met\'alicas conectadas a una diferencia de potencial; cuando est\'an 
sumergidas en agua, se registra un voltaje $V_w$ al que le asociamos el 100\% de humedad y luego, cuando las placas est\'an al aire, se registra un voltaje $V_a$ al que le 
asociamos el 0\% de humedad; por lo tanto hay 2 puntos $(V_w,100\%)$ y $(V_a,0\%)$ con los cuales se hace un ajuste lineal (una recta que pase por los 
dos) y se obtiene:
\begin{equation}
	H(x)=\frac{100}{V_{w}-V_{a}}(V-V_{a})
\end{equation}
donde $H$ es la humedad en la tierra y $V$ es el voltaje registrado por el sensor. Este driver era utilizado para el sensor FC-28 de arduino, ahora se utiliza un sensor 
similar creado por Diego Cabrer. Con el nuevo sensor, el driver responde correctamente al 
0\% y al 100\%; para los valores intermedios arroja lecturas l\'ogicas, sin embargo ser\'ia interesante calibrar 
el sensor con 9 muestras de tierra al $10\%, 20\%, 30\%, 40\%, 50\%, 60\%, 70\%, 80\%, 90\%$ de humedad y generar 11 puntos experimentales

\begin{equation}
(V_{a},0\%) , (V_{10\%},10\%), ... , (V_{w},100\%)
\end{equation}

para despu\'es hacer un ajuste (no necesariamente lineal) y obtener una mejor funci\'on $H(V)$.

\section{Base de datos}
La base de datos sirve para registrar toda la informaci\'on de los 
sistemas: los actuadores y sensores, el tiempo y periodo de 
riego, la temperatura del agua, el rango de temperatura y humedad 
ambiente, entre otras cosas. La base tambi\'en registra las lecturas de 
los sensores y las acciones que efect\'uan los actuadores. De las 
lecturas se procesan los promedios por hora y se 
guardan en una tabla de estad\'isticas.
 
\subsection{Informaci\'on del sistema} 
Para la informaci\'on de los sistemas se tienen varias tablas. La tabla 
\verb|sistema| guarda la informaci\'on del sistema: nombre, cultivo, id, url, ubicaci\'on y si es sistema principal o subsistema. La 
tabla \verb|sensor| tiene la informaci\'on de los sensores: el tipo de sensor, 
el pin del microprocesador al que est\'a conectado el sensor, 
el periodo de muestreo y el identificador del sistema al que pertenece. La tabla \verb|actuador| tiene la informaci\'on de los actuadores: el pin al que est\'a conectado el 
actuador, 
el tipo de actuador y el identificador del sistema al que pertenece. La tabla \verb|riego| contiene la informaci\'on del
 riego: el tiempo de riego, el periodo de riego y la temperatura del agua con la que se riega. Existen 3 valores de tiempo, 3 valores para el periodo y 3 
valores de temperatura del agua; cada uno est\'a relacionado con la temperatura ambiente (la relaci\'on se explicar\'a m\'as adelante). La tabla 
\verb|planificacion| contiene 
la informaci\'on del cultivo: nombre del cultivo, rango ideal de 
temperatura y humedad ambiente, y rango de ph.

\subsection{Monitoreo}
Las tablas de monitoreos son 

\begin{verbatim}
monitoreo_actuador, monitoreo_sensor
\end{verbatim}

en la tabla \verb|monitoreo_sensor| se registran fecha, hora, valor de la lectura e id del sensor; 
en la tabla \verb|monitoreo actuador| se registran fecha, hora, tipo de acci\'on y valor 
de activaci\'on. El campo tipo acci\'on de monitoreo actuador se 
refiere a encendido (1) o apagado (0).

\subsection{Estad\'isticas}
La tabla \verb|estadistica_hora| se llena con los promedios por hora de los 
registros \verb|monitoreo_sensor|. Los campos de esta tabla corresponden a la fecha, el tipo de sensor, el id sistema del cual se hacen los promedios y las horas del d\'ia. 
Dentro de los campos de las horas del d\'ia, se guardar\'a el promedio de las lecturas que se hayan efectuado a esa hora, por ejemplo, un registro quedar\'ia de la siguiente 
manera:  
\begin{tabular}[h]{|l|l|l|l|l|l|}
        \hline
        fecha & tipo & sistema & 00 & ... & 23 \\ \hline
        2018-05-15 & ph & 3 & 7.5 & ... & 8.1 \\  \hline
\end{tabular}

Este registro indica el promedio de ph del sistema con id$=3$ a las 00 horas, 01 horas y as\'i sucesivamente hasta mostrar el promedio de ph a las 23 horas. La tabla se 
actualiza mediante un trigger que se activa al introducirse un nuevo registro en la tabla \verb|monitoreo_sensor|. 

\subsection{Conector c++ con mysql}
Las consultas a la base de datos desde el planificador se hacen con \verb|Mysql/C++ connector|, es un driver que provee Mysql. El driver se instala con el paquete 
\verb|libmysqlcppconn-dev|
en la parte de configuraci\'on. Para ejecutar 
consultas en paralelo, se cre\'o una clase llamada \verb|conector|, definida en el archivo \verb|Codes/conector.cpp| y declarada en \verb|Codes/Headers/conector.cpp|. 
Las consultas dentro del programa se ejecutan de la siguiente manera:

\begin{verbatim}
conector C;
C.conectar();
//para ejecutar sin resultado
C.stmt->execute("consulta");
//con resultado
C.res=C.stmt->executeQuery("consulta");
//...hacer cosas
C.borrar();
\end{verbatim}


\section{Herramientas del sistema}
Las herramientas del sistema son un conjunto de estructuras y funciones usadas por el planificador de tareas y las pruebas del hardware; se 
encuentran declaradas y definidas en los 
archivos \verb|Codes/Headers/tools.cpp| y \verb|Codes/tools.cpp| respectivamente. Las estructuras m\'as importantes son: \verb|sensor|, 
\verb|actuador| e \verb|hilo|; fueron definidas de la siguiente manera:

\begin{verbatim}
struct sensor{
        int id,muestreo,pin,sys;
        float valor;
        string name;
        };
struct actuador{
        int id,pin,seg,sys;
        string name;
        };
struct hilo{
        string name;
        pid_t tid;
        };
\end{verbatim}
Las variables en la estructura \verb|sensor| corresponden a la siguiente informaci\'on: 

\begin{itemize}
  \item \verb|id|: identificador del sensor, es un n\'umero entero.
  \item \verb|muestreo|: minutos entre cada muestreo (periodo de muestreo). 
  \item \verb|pin|: pin al que est\'a conectado el sensor.
  \item \verb|sys|: identificador del sistema al que pertenece el sensor.
  \item \verb|valor|: valor de la variable del entorno que representa el sensor. 
  \item \verb|name|: tipo de sensor, sirve para relacionar el sensor con el driver. 
\end{itemize}

La estructura \verb|actuador| 
es an\'aloga a \verb|sensor| para las variables de \verb|id|, \verb|pin|, \verb|sys| y \verb|name|. La variable diferente \verb|seg|, representa los segundos de acci\'on 
del actuador. La estructura \verb|hilo| sirve para almacenar la informaci\'on de un thread; la variable \verb|name| guarda el nombre del thread y la variable \verb|tid| 
(thread id) guarda el identificador del thread.

\section{Pruebas hardware}
Para verificar el funcionamiento de los sensores y actuadores, se cre\'o el ejecutable \verb|Hardware_test| definido en el archivo \verb|Codes/pruebas.cpp|. El ejecutable 
carga 
los sensores y actuadores registrados en la base de datos dentro de las estructuras definidas anteriormente. Una vez creados los objetos con la informaci\'on, se muestra un 
men\'u donde se decide si se desea probar los actuadores o los sensores. 
Posteriormente, se muestra un men\'u y se debe escoger si emitir una prueba general o escoger el sensor/actuador que se desee probar. La interfaz gr\'afica es muy 
rudimentaria por lo cual ser\'ia buena idea trabajar en este aspecto del proyecto. El ejecutable no tiene manera de hacer m\'ultiples pruebas o de suspender la ejecuci\'on 
de una 
prueba; la soluci\'on ser\'ia ejecutar en threads las pruebas y 
regresar al men\'u para decidir si ejecutar otra prueba o terminar alguna.

\section{Planificador de tareas}
El planificador de tareas es un ejecutable llamado \verb|Hidroponia| creado del c\'odigo \verb|scheduller.cpp| en la carpeta Codes. El ejecutable tiene un par\'ametro 
entero que representa el identificador del sistema que se va a activar. Por ejemplo, dentro de la carpeta \verb|Hidroponia|, ejecutando el comando:

\begin{verbatim}
./Hidroponia 2
\end{verbatim}

se inicia el planificador de tareas para el sistema con identificador 2 en la tabla \verb|sistema| en la base de datos. Las tareas del sistema son las siguientes:

\begin{enumerate}  
	\item Capturar informaci\'on de la base 
	\item Programar rutinas de los sensores
	\item Programar el riego y los actuadores
	\item Monitoreo de sensores y actuadores
	\item Gestionar la ca\'ida de alg\'un thread
\end{enumerate}

\subsection{Captura de datos}
La captura de datos comienza por la informaci\'on del sistema; partiendo del id se consulta lo que hay en la tabla \verb|sistema| y se guarda en una estructura 
llamada \verb|sistema|. Despu\'es se 
consulta su planificador y el riego para almacenar los datos en una clase llamada 
\verb|planificador|, declarada en \verb|tools.hpp| y definida en \verb|tools.cpp|. Posteriormente, se cuenta la cantidad de sensores y actuadores para crear dos arreglos de 
estructuras \verb|sensor| y \verb|actuador| de tama\~nos respectivos; creados los arreglos, se inicializan con la informaci\'on de cada sensor del sistema. Esta fase del 
programa consiste en guardar la informaci\'on de la base de datos en las estructuras definidas en las herramientas del sistema (tools.cpp). 
 
\subsection{Rutinas de los sensores}
Lo siguiente es ejecutar las rutinas de muestreo; para cada sensor se lanza un thread que realiza dicha rutina. El pseudoc\'odigo del thread de muestreo es el siguiente:

\begin{verbatim}
Muestreo(sensor){
  loop{
    esperar el pin;
    ocupar pin;
    inicializar sensor.valor=-1;
    intentos=0;
    while(sensor.valor==-1){
      indexar el sensor con su driver
      if(sensor.name=tipo1){
        sensor.valor=driver_tipo1(sensor.pin);
        }
      else if(sensor.name=tipo2){
        sensor.valor=driver_tipo2(sensor.pin);
        } 
      ...
      ...
      else{
        no hay driver;
        exit;
        }
      if(sensor.valor==-1&intentos++<2){
        reportar error;
        }
      }fin while
    registrar lectura
    en la base de datos;
    liberar pin;
    sleep(sensor.muestreo);
    }fin loop
}fin muestreo
\end{verbatim}

La edici\'on de este thread es la 
\'ultima parte 
de la automatizaci\'on del protocolo para a\~nadir drivers; se debe colocar la funci\'on en un archivo separado para editarlo, actualmente se encuentra en 
\verb|Codes/scheduller.cpp| junto con todo el planificador. Para continuar con el protocolo, se tendr\'ian que a\~nadir las siguientes l\'ineas de c\'odigo despu\'es 
del \'ultimo \verb|else if| dentro del thread:

\begin{verbatim}
else if(sensor.name=="tipo_nuevo"){
  sensor.valor=nuevo_driver(sensor.pin);
  }
\end{verbatim}

despu\'es compilar la funci\'on en un objeto y actualizar las librer\'ias en el archivo \verb|libs.hpp|. Finalmente, actualizar los ejecutables \verb|Hardare_test| e 
\verb|Hidroponia| para tener las pruebas del sistema y el planificador de tareas actualizados con el nuevo driver.

\subsection{Actuadores}
Los diferentes tipos de actuadores que soporta el sistema son:
\begin{enumerate}
	\item bomba de agua
	\item calefactor del agua
	\item ventilador
	\item reguladores de ph
\end{enumerate}
y para cada uno est\'a definido un thread que se lanza s\'olo si se cubren determinados requerimentos. Para el riego existen 2 threads: uno se ejecuta 
peri\'odicamente 
de acuerdo a la informaci\'on del riego en la 
base de datos y el otro se activa dependiendo la humedad de la tierra (si baja a determinado nivel se activa el riego hasta que la humedad vuelva a subir). El primer thread 
de riego s\'olo 
requiere una bomba de agua mientras que el otro requiere una bomba de agua y un sensor de humedad en tierra para poder ejecutarse. El thread del calefactor de agua requiere 
un actuador de tipo calefactor y un sensor de temperatura en liquidos; el principio es el mismo que el segundo riego, si el agua baja a cierta temperatura se activa el 
calefactor 
hasta que se consiga una temperatura en el rango asignado por el planificador. El ventilador es an\'alogo, se activa a determinada temperatura ambiente y se desactiva 
cuando la temperatura ambiente llega al rango ideal. Los reguladores de ph no se instalaron pero el thread debe ser creado de la misma manera que los otros 
actuadores; para subir el ph, el pseudoc\'odigo ser\'ia:

\begin{verbatim}
activar(regulador_phmas.pin);
while(sensorph.valor<phminimo){
  sleep(t);
  sensorph.valor=lectura de ph;
  }
desactivar(regulador_phmas.pin3);
\end{verbatim}

Para disminuir el ph, el pseudoc\'odigo ser\'ia igual \'unicamente se modificar\'ia la condici\'on del while: \verb|while(sensorph.valor>phmax)|. La letra \verb|t| 
representa un peque\~no periodo de muestreo para saber cu\'ando desactivar 
el actuador correspondiente.

\subsection{Sem\'aforos para los threads de muestreos}
Dos de los problemas relacionados con el multithreading fueron deadlocks al momento de hacer consultas a la base y la concurrencia de threads que intentaban acceder al 
mismo espacio de memoria en un determinado instante; dichos problemas causaban el fallo del sistema operativo. El primer problema se resolvi\'o con un 
thread (sem\'aforo)
encargado de ejecutar las consultas de monitoreo en orden. El sem\'aforo de los monitoreos consiste en una fila donde se insertan las 
consultas en forma de strings y se van 
ejecutando una por una. Para el segundo problema, el sem\'aforo consiste en un arreglo de valores booleanos y 3 pasos: esperar, usar y liberar. Cada elemento del 
arreglo 
corresponde a un pin del duino; si el elemento es verdadero quiere decir que el pin correspondiente al index del elemento est\'a ocupado. Cuando un thread quiere usar 
el pin, debe consultar el arreglo para ver si est\'a ocupado, de ser as\'i debe esperar. La forma de usar un pin es la siguiente:

\begin{verbatim}
while(Pinbussy[4]){ //esperar }
//establecer el pin ocupado
Pinbussy[4]=true;

...usar el pin

//liberar el pin
Pinbussy[4]=false;
\end{verbatim}

En este ejemplo, el thread espera en el loop \verb|while| hasta que se libera el pin 4; saliendo del loop, se fija el valor en \verb|true| para establecerlo como ocupado; 
posteriormente, se utiliza el pin y se libera. Los sem\'aforos se encuentran 
definidos en un thread llamado \verb|monitoreo| y en un conjunto de funciones llamadas \verb|usar(pin)|, \verb|liberar(pin)| y \verb|esperar(pin)|. El \'ultimo sem\'aforo 
creado fue el administrador de sensores de temperatura en l\'iquidos. El driver del sensor de temperatura en l\'iquidos utiliza un archivo para su funcionamiento; cuando 
dos threads intentaban acceder al archivo al mismo tiempo, el sistema operativo colapsaba. El sem\'aforo es un programa que consulta los
sensores de temperatura en l\'iquidos dados de alta para despu\'es efect\'uar lecturas con cada uno y posteriormente registrarlas 
en la base de datos. Las muestras y los registros se efect\'uan de manera ordenada, un sensor a la vez, para evitar la convergencia de threads en el archivo de lecturas.

\subsection{Recuperador de threads}
Otro problema que surgi\'o con el multithreading fue la muerte del thread de monitoreo. A pesar del thread ca\'ido, el programa continu\'o en ejecuci\'on; no se registraban 
los muestreos a la base de datos monitoreo. En aquel momento no hubo manera de rastrear el error que ocacion\'o el fallo del thread. Para poder rastear los errores y las 
ca\'idas del thread fue necesario crear archivos log; uno para el registro de los procesos normales y otro para los errores. Los archivos log se crearon 
redirigiendo las variables \verb|stdout| y \verb|stderr| hacia los archivos \verb|Hidroponia/Logs/logfile| y \verb|Hidroponia/Logs/errfile| respectivamente. Cada 
planificador de tareas tiene su \verb|logfile| y \verb|errfile|, se encuentran diferenciados por el identificador del sistema al final del nombre del archivo. Para manejar 
la 
ca\'ida de threads se cre\'o un 
reparador de threads dentro del loop principal del administrador de tareas. El pseudoc\'odigo es el siguiente:

\begin{verbatim}

main(int id_planificador){
...cargar sensores y actuadores
n=no. de threads;
(usamos la estructura hilo)
arreglo de hilos[n];
...Ejecutar threads y guardar info 
...en estructura hilo[i]
loop{
  for(i=0;i<n;i++){
    if(!running(hilos[i].tid){
      relanzar el thread;
      reportar en logfile;
      }end if
    }end for
  esperar media hora;
  }end loop
}end main
\end{verbatim}

En la rutina se verifica peri\'odicamente si cae alg\'un thread y, en caso de caerse, relanzarlo y reportar. La funci\'on \verb|running(hilo.tid)| verifica si est\'a 
corriendo un thread mediante su tid; la funci\'on se apoya de una carpeta que contiene archivos con informaci\'on de status para cada thread que sea ejecutado; la carpeta es 
\verb|Hidroponia/Threads|. 
El reparador de threads 
es el loop principal del planificador de 
tareas, mantiene el programa en ejecuci\'on; si llegara a caer el reparador de threads, caer\'ian todos los threads con \'el. Hace falta manejar la ca\'ida del thread 
principal, debe ser programado en el sistema operativo; alg\'un script que est\'e ejecutando peri\'odicamente la opci\'on 2 del administrador de sistemas.

\section{Generador de estad\'isticas}
El generador de estad\'isticas es un ejecutable llamado \verb|stats| compilado del c\'odigo \verb|stats.cpp|. Lo que hace es generar peri\'odicamente los 
promedios de las lecturas y registrarlos en la tabla \verb|estadistica_hora|. El pseudoc\'odigo del 
generador de estad\'isticas es el siguiente:
\begin{verbatim}
loop{
  consultar los sistemas activos;
  para cada sistema{
    consultar sensores registrados;
    para cada sensor{
      if(registros_sensor==0){
        crear registro;
        }
      avg=promedio de lecturas 
          en la hora actual;
  update 
   estadistica_hora.hora_actual=avg;
      }fin sensores
    }fin sistemas
  esperar 20 minutos;
  }fin loop;
\end{verbatim}

El programa fue reemplazado por el trigger que promedia en la hora con cada nuevo registro en la tabla \verb|monitoreo_sensor|.

\section{Administrador de sistemas}
El adminstrador de sistemas es un ejecutable llamado \verb|admin| definido en el c\'odigo \verb|manager.cpp| en la carpeta Codes. El administrador tiene 4 opciones, 3 de 
ellas funcionales y 1 en desarrollo. Las opciones se ejecutan sobre la carpeta principal \verb|Hidroponia| de la siguiente manera:

\begin{verbatim}
./admin opc
\end{verbatim}

donde \verb|opc| es un entero que representa la opci\'on que se desea ejecutar. Las opciones son: ''2'' para actualizar los sistemas con la base de datos, ''3'' para 
actualizar 
los 
drivers, ''4'' matar todo y ''5'' ejecutar las pruebas del hardware.

\subsection{Opci\'on 2}
La opci\'on 2 actualiza los ''schedullers'' en ejecuci\'on con la base de datos. El pseudoc\'odigo es el siguiente:

\begin{verbatim}
load_systems(){
  get sistemas where estado=1;
  for each sistema{
    if(!running(sistema)){
      ./Hidroponia sistema.id;
      get subsistemas where estado=1;
      for each subsistema{
        if(!running(subsistema)){
          ./Hidroponia subsistema.id;
          }end if
        }end subsistemas
      }end if
    }end sistemas
  get sistemas with estado=0;
  for each sistema{
    if(running(sistema)){
      kill(sistema);
      get subsistemas;
      for each subsistema{
        kill(subsistema);
        }end subsistemas
      }end if
    }end sistemas
  if(!running(stats)){ ./stats }
  if(!running(wateradmin)){ ./wateradmin }
  }end load_systems
\end{verbatim}

Por lo tanto, si se quisiera apagar o encender un determinado sistema ser\'ia 
recomendable editar el campo \verb|estado| en la base de datos y ejecutar la opci\'on 2 del administrador de sistemas.

\subsection{Opci\'on 3}
Esta opci\'on sigue en desarrollo y se ha comentado a lo largo del documento. Sirve para automatizar el protocolo para a\~nadir un driver al sistema. Lo que se tiene hasta 
ahora es la creaci\'on de los objetos con los archivos de la carpeta Drivers y la creaci\'on del archivo \verb|libs.hpp| que 
contiene todos los headers; se logra con la ejecuci\'on del binario \verb|comp| sobre la carpeta \verb|Hidroponia|. Falta editar el c\'odigo del thread de muestreo y 
actualizar los binarios \verb|Hidroponia| y 
\verb|Hardware_test|.

\subsection{Opci\'on 4}
Esta opci\'on consulta todos los sistemas registrados en la base de datos y apaga cada uno de ellos. Tambi\'en termina los programas \verb|stats| y \verb|wateradmin|. 
Se utiliza para reiniciar los sistemas, primero ejecutar opci\'on 4 y luego opci\'on 2.

\subsection{Opci\'on 5}
Es la opci\'on para ejecutar las pruebas del hardware. Esta opci\'on primero termina todos los sistemas con la opci\'on 4 y despu\'es ejecuta el 
binario \verb|./Hardware_test| para correr las pruebas del hardware.

\section{Makefile}
El archivo \verb|makefile|, ubicado en la carpeta principal, contiene las definiciones de los comandos de compilaci\'on para las partes del proyecto. Al ejecutar el comando 
\verb|make|, se compilan todos los m\'odulos del sistema hidrop\'onico. Para compilar los m\'odulos por separado, se ejecuta el comando \verb|make modulo| donde los 
diferentes m\'odulos son:
\begin{itemize}
\item \verb|compiler|: compila el c\'odigo \verb|compiler.cpp| para hacer el ejecutable \verb|comp|.
\item \verb|drivers|: ejecuta \verb|./comp| para crear los objetos.
\item \verb|conector|: compila el archivo \verb|conector.cpp| para la clase conector.
\item \verb|tools|: compila \verb|tools.cpp|.
\item \verb|pruebas|: compila \verb|pruebas.cpp| para crear \verb|Hardware_test|.
\item \verb|hidroponia|: compila \verb|scheduller.cpp| para crear el ejecutable \verb|Hidroponia|.
\item \verb|manager|: compila \verb|manager.cpp| para crear \verb|admin|.
\item \verb|statistics|: compila \verb|stats.cpp| para crear \verb|stats|.
\item \verb|water|: compila \verb|wateradmin.cpp| para crear \verb|wateradmin|.
\item \verb|clean|: borra los ejecutables.
\end{itemize}

Por ejemplo, para actualizar el planificador de tareas, despu\'es de una modificaci\'on al c\'odigo, se debe ejecutar \verb|make hidroponia| sobre la carpeta principal 
\verb|Hidroponia|.

\section{Interfaz web}
Se cre\'o una peque\~na interfaz gr\'afica provicional para agregar/editar sistemas, planificadores y riegos; tambi\'en  permite monitorear las lecturas de los sensores y 
las acciones de 
los actuadores as\'i como las estad\'isticas generadas en la tabla \verb|estadistica_hora|. 
Los c\'odigos de la interfaz se encuentran en la carpeta

\begin{verbatim}
/var/www/Proyecto/
\end{verbatim}

Para accesar a la interfaz desde el navegador, se debe ingresar en la barra de direcci\'on lo siguiente:

\begin{verbatim}
10.10.254.78/Proyecto
\end{verbatim}

Es importante resaltar que ser\'ia la ip que tuviera asignada el servidor.

\subsection{Sistemas}
En esta parte de la interfaz se registra/edita un sistema o subsistema. Al pasar a la p\'agina \verb|sistema.php| se muestra la opci\'on de editar un sistema o 
crear uno nuevo. Para la opci\'on de editar se despliegan todos los sistemas registrados; si se desea editar un subsistema hay que escoger el sistema padre y despu\'es 
el subsistema. En la parte de edici\'on del sistema o subsistema se encuentran las siguientes opciones:

\begin{itemize}
\item A\~nadir, quitar o modificar sensores y actuadores
\item Editar la informaci\'on del sistema
\item Cambiar de planificador o riego
\item Encender el sistema
\item Borrar el sistema
\end{itemize}


\subsection{Planificaciones y Riegos}
An\'alogo a sistemas, se muestran las opciones de crear o de editar planificaciones y riegos Los c\'odigos de las p\'aginas son \verb|planificador.php| y \verb|riego.php| 
respectivamente.

\subsection{Monitoreo}
Se encuentra en el archivo \verb|monitoreo.php|; lo que hace el script es filtrar las tablas \verb|monitoreo_sensor| y \verb|monitoreo_actuador| pregunt\'andole al 
usuario el sistema que desea monitorear, si desea monitorear sensores o actuadores, qu\'e tipo de sensor o actuador y el intervalo que desea monitorear. Una vez 
que se fijan esas condiciones, se genera una consulta a la base de datos y se crean p\'aginas de 30 registros para visualizar las tablas de monitoreo filtradas. 
Para mayor detalle, consultar el c\'odigo \verb|monitoreo.php| en la carpeta de la interfaz.

\subsection{Estad\'isticas}
En esta parte de la interfaz se pueden visualizar los promedios de las variables de entorno respecto cada hora, d\'ia, semana o mes. Para la visualizaci\'on de promedios 
por hora se hace una consulta a la base de datos de la tabla \verb|estadistica_hora| donde se guardan y se generan los promedios de las lecturas que est\'an en 
\verb|monitoreo_sensor| de manera ordenada en registros diarios. Para la visualizaci\'on en d\'ias se consulta a la tabla \verb|monitoreo_sensor| el promedio de los 
valores de cierto sensor a lo largo del d\'ia, an\'alogamente para semana y mes. Para mayor detalle de las consultas, revisar el c\'odigo \verb|estadisticas.php| 
en 
la carpeta de la interfaz.

\end{document}
